\documentclass{article}

\usepackage{graphicx}
\usepackage{tikz}
\usepackage{hyperref}
\usepackage[margin=1in]{geometry}
\usepackage{enumitem}
\usetikzlibrary{arrows.meta, positioning, calc}

\begin{document}

\title{Updates on data processing pipeline}
\author{}
\date{09 January 2026 \ \small Reference: 23 December 2025}
\maketitle

\section*{Key Changes}
\begin{itemize}[leftmargin=*]
    \item Patch extraction was rebuilt: positives come from the malignant mask or its bounding box with a minimum lesion coverage, negatives are split into near/far groups with coverage caps, intensity filters, and balanced 1:1 targets for malignant vs.\ healthy patches.
    \item Patient-level train/val/test splitting is now supported (seeded) to avoid leakage, with per-split patch folders and patient IDs carried through loading and feature building.
    \item Config, selection, and models with support for the new sampling features (coverage thresholds, near/far radii, nonzero/intensity filters) plus group subsampling and GroupShuffleSplit based validation; mean/std are only stored when standardizing and AUC handling is more robust.
    \item Pipeline now loads arbitrary splits, to keep patient labels through response computation, evaluates trained models on additional splits, and logs per split metrics alongside ROC/PR and confusion plots.
    \item Drop patches containing dark pixels more than 35\%.
\end{itemize}

\begin{figure}[ht]
    \centering
    \begin{tikzpicture}

        % Left images
        \node (L1) {\includegraphics[width=0.22\linewidth]{23_december_images/healthy_old.png}};
        \node (L2) [below=0.2cm of L1] {\includegraphics[width=0.22\linewidth]{23_december_images/cancer_old.png}};

        % Right images
        \node (R1) [right=3cm of L1] {\includegraphics[width=0.22\linewidth]{23_december_images/healthy_new.png}};
        \node (R2) [below=0.2cm of R1] {\includegraphics[width=0.22\linewidth]{23_december_images/cancer_new.png}};

        % Midpoints of left and right columns
        \coordinate (Lmid) at ($(L1)!0.5!(L2)$);
        \coordinate (Rmid) at ($(R1)!0.5!(R2)$);

        % Arrow with centered text
        \draw[-{Stealth}, thick]
            (Lmid) -- node[midway, above]{\small Old to New} (Rmid);

    \end{tikzpicture}
    \caption{Inspection of data.}
    \label{fig:new_data_pipeline}
\end{figure}

\newpage

I was expecting more changes in decision boundary graph but I got a similar pattern. I ran multiple experiments with different configuration of kernel selection (ABESS/MMR, Number of Kernels, K best kernels) but unfortunately I could not get a better distinction.

\begin{figure}[ht]
    \centering
    \begin{tikzpicture}

        % Left images
        \node (L1) {\includegraphics[width=0.22\linewidth]{23_december_images/scatter_best_pair_exp5.png}};

        % Right images
        \node (R1) [right=3cm of L1] {\includegraphics[width=0.22\linewidth]{23_december_images/scatter_best_pair_exp7.png}};

        \draw[-{Stealth}, thick]
            (L1.east) -- node[above]{\small New Data Pipeline} (R1.west);

    \end{tikzpicture}
    \caption{Overview of the new data pipeline from original samples to processed outputs.}
    \label{fig:new_data_pipeline}
\end{figure}

\section*{More changes and a new experiment}
\begin{itemize}[leftmargin=*]
    \item In the previous experiments, I was using the non malignant parts of the image (healthy parts of malignant image) as healthy, in the following I included patches from truly healthy data as well.
    \item Feature side: per-patch responses are z-scored before convolution and default aggregation switched to \texttt{abs\_max}. (I think not a good idea)
    \item Standardize patches using their mean and std.
\end{itemize}
I managed to make a less reliable experiment using these 3 changes:

\begin{figure}[ht]
    \centering
    \begin{tikzpicture}

        % Left images
        \node (L1) {\includegraphics[width=0.3\linewidth]{23_december_images/scatter_best_pair_exp9.png}};

        % Right images
        \node (R1) [right=1cm of L1] {\includegraphics[width=0.3\linewidth]{23_december_images/scatter_best_triple_exp9.png}};

    \end{tikzpicture}
    \caption{}
    \label{fig:new_data_pipeline}
\end{figure}

\section*{Pairwise preview}
\begin{figure}[ht]
    \centering
    \includegraphics[width=0.55\linewidth]{../results/pairwise_preview/selected/pairwise_correlation_heatmap.png}
    \caption{Selected pairwise correlation heatmap from the latest run.}
    \label{fig:pairwise_heatmap}
\end{figure}

\begin{figure}[p]
    \centering
    \includegraphics[width=0.23\linewidth]{../results/pairwise_preview/selected/scatter_k236_k314.png} \hfill     \includegraphics[width=0.23\linewidth]{../results/pairwise_preview/selected/scatter_k236_k601.png} \hfill     \includegraphics[width=0.23\linewidth]{../results/pairwise_preview/selected/scatter_k236_k810.png} \hfill     \includegraphics[width=0.23\linewidth]{../results/pairwise_preview/selected/scatter_k236_k828.png}
    \\[0.2cm]
    \includegraphics[width=0.23\linewidth]{../results/pairwise_preview/selected/scatter_k810_k601.png} \hfill     \includegraphics[width=0.23\linewidth]{../results/pairwise_preview/selected/scatter_k828_k631.png} \hfill     \includegraphics[width=0.23\linewidth]{../results/pairwise_preview/selected/scatter_k828_k810.png} \hfill     \includegraphics[width=0.23\linewidth]{../results/pairwise_preview/selected/scatter_k869_k236.png}
    \caption{Selected pairwise scatter previews (1/2)}
    \label{fig:pairwise_scatter_1}
\end{figure}

\begin{figure}[p]
    \centering
    \includegraphics[width=0.23\linewidth]{../results/pairwise_preview/selected/scatter_k869_k828.png} \hfill     \includegraphics[width=0.23\linewidth]{../results/pairwise_preview/selected/scatter_k905_k236.png} \hfill     \includegraphics[width=0.23\linewidth]{../results/pairwise_preview/selected/scatter_k905_k314.png} \hfill     \includegraphics[width=0.23\linewidth]{../results/pairwise_preview/selected/scatter_k905_k631.png}
    \\[0.2cm]
    \includegraphics[width=0.23\linewidth]{../results/pairwise_preview/selected/scatter_k929_k236.png} \hfill     \includegraphics[width=0.23\linewidth]{../results/pairwise_preview/selected/scatter_k929_k314.png} \hfill     \includegraphics[width=0.23\linewidth]{../results/pairwise_preview/selected/scatter_k929_k828.png} \hfill     \includegraphics[width=0.23\linewidth]{../results/pairwise_preview/selected/scatter_k929_k869.png}
    \caption{Selected pairwise scatter previews (2/2)}
    \label{fig:pairwise_scatter_2}
\end{figure}

\noindent\textbf{Short notes.} The heatmap summarizes kernel-response similarity to flag redundancy, and the selected scatter grids show class overlap patterns for the curated kernel pairs.

Current plans:
\begin{itemize}
    \item Rebalance negatives: reduce/disable malignant-background negatives and keep most negatives from healthy slides; also increase far neg radius.
    \item Loosen positive strictness slightly: drop min\_pos\_coverage back toward 0.6–0.7; keep bbox positives off (bbox around the lesion tissue).
    \item Try response\_fn="abs\_max" vs mean\_abs and turn off patch z-scoring.
    \item Model: standardize features, use logistic (not MLP) for the 2D/3D plots to avoid overfitting on small subsets; rerun subset search with a fixed train\_subset\_frac (e.g., 0.5) for stability.
    \item And a lot more experiments...
\end{itemize}
\href{https://github.com/Masoudvahid/Kernel-Based-Cancer-Classification}{link to github}


\end{document}
